\subsubsection{Eine Aufbauphase als gewöhnliche Funktion}

Ein Paar \((u,x)\) hält zunächst nur fest, dass dem Wort \(u\) der
Wert \(x\) zugeordnet wird. Alle solchen Paare liegen in
\(B=W\times X\). Zu einer vorliegenden Zuordnung \(H\subseteq B\)
bilden wir das Anfangspaar und sämtliche Fortsetzungen:
\[
  T(H)=\{(e,x_0)\}\cup
  \{(s(u,a),r(x,a))\mid (u,x)\in H,\ a\in A\}.
\]
Dies ist eine gewöhnliche Mengenbildung aus den gegebenen Funktionen
\(s,r\). Insbesondere wird auf der rechten Seite kein gesuchter
Rekursionswert vorausgesetzt. Die folgenden Definitionen schreiben diese
Mengenbildung mit der Aussonderung und der Termgraphkonstruktion aus
Band~5 aus.

\Needspace{14\baselineskip}
\FormulaDefDeltaK[Umgebungsmenge der Wort-Wert-Paare]{%
  B\coloneqq W\times X
}{WordRecursionAmbientDef}{\DeltaRow{Mengen}{W\dsep X}
  \DeltaRow{Neues Symbol}{B}}

\Needspace{14\baselineskip}
\FormulaDefDeltaK[Fortsetzung einer vorliegenden Zuordnung]{%
  \begin{aligned}[t]
    \tau(H)\coloneqq\{(e,x_0)\}\cup\{p\in B\mid
      &\exists u\in W\,\exists x\in X\,\exists a\in A\;\\[-2pt]
      &((u,x)\in H\land p=(s(u,a),r(x,a)))\}
  \end{aligned}
}{WordRecursionStageTermDef}{%
  \DeltaRow{Mengen}{A\dsep W\dsep X\dsep e\dsep x_0\dsep H\dsep p\dsep u\dsep x\dsep a}
  \DeltaRow{Funktionen}{s\dsep r}\DeltaRow{Neues Symbol}{\tau(H)}}

\Needspace{14\baselineskip}
\FormulaDefDeltaK[Die Funktion für eine Aufbauphase]{%
  T\coloneqq\{(H,K)\in\powerset(B)\times\powerset(B)\mid K=\tau(H)\}
}{WordRecursionStageOperatorDef}{%
  \DeltaRow{Mengen}{H\dsep K}\DeltaRow{Neues Symbol}{T}}

Die Hilfsschreibweise \(\tau(H)\) bezeichnet den Mengenterm;
\(T\) ist dessen Funktionsgraph. Der nächste Nachweis stellt sicher,
dass wir \(T\) als Schrittfunktion einer Zahlenrekursion verwenden
dürfen. Die Beschränkung des Schrittbilds auf \(B\) verliert keine
Fortsetzung: W0 und die Typisierung von \(r\) liefern gerade die
benötigte Zugehörigkeit zu \(W\times X\).

\Needspace{14\baselineskip}
\FormulaThmDeltaK[Die Aufbauphase ist eine Funktion auf der Potenzmenge]{%
  \begin{aligned}[t]
    &\mathsf{WortStr}_A(W,e,s)\dsep x_0\in X\vdash{}\\[-2pt]
    &T\colon\powerset(B)\to\powerset(B)\land
      \forall H\in\powerset(B)\;T(H)=\tau(H)
  \end{aligned}
}{WordRecursionStageOperatorFunction}{%
  \DeltaRow{Mengen}{A\dsep W\dsep e\dsep X\dsep x_0\dsep H}\DeltaRow{Funktionen}{s\dsep r}}
\begin{tabproofwide}
  \proofstepwidestar[1]{\mathsf{WortStr}_A(W,e,s)}{\rA}
  \proofstepwidestar[2]{x_0\in X}{\rA}
  \proofstepwidestar[1]{e\in W}{\rAEa{\FormulaRefAuto{WordStructureTypingAxiom}[ax]{1}}}
  \proofstepwidestar[1,2]{(e,x_0)\in B}{\rIE{\FormulaRefAuto{a=b\vdash b=a}[thm]{\FormulaRefAuto{WordRecursionAmbientDef}[def]},\FormulaRefAuto{(a,b)\in A\times B\eqvdash a\in A\land b\in B}[thm]{\rAI{3,2}}}}
  \proofstepwidestar[1,2]{\{(e,x_0)\}\subseteq B}{\FormulaRefAuto{a\in A\vdash\{a\}\subseteq A}[thm]{4}}
  \proofstepwidestar[]{\{p\in B\mid\exists u\in W\,\exists x\in X\,\exists a\in A\;((u,x)\in H\land p=(s(u,a),r(x,a)))\}\subseteq B}{\FormulaRefAuto{\{x\in A\mid P(x)\}\subseteq A}[thm]}
  \proofstepwidestar[1,2]{\tau(H)\subseteq B}{\rIE{\FormulaRefAuto{a=b\vdash b=a}[thm]{\FormulaRefAuto{WordRecursionStageTermDef}[def]},\FormulaRefAuto{A\subseteq C\dsep B\subseteq C\vdash A\cup B\subseteq C}[thm]{5,6}}}
  \proofstepwidestar[1,2]{\tau(H)\in\powerset(B)}{\FormulaRefAuto{x\in\powerset(A)\eqvdash x\subseteq A}[thm]{7}}
  \proofstepwidestar[9]{H\in\powerset(B)}{\rA}
  \proofstepwidestar[1,2]{\forall H\in\powerset(B)\;\tau(H)\in\powerset(B)}{\rUI{\rRI{9,8}}}
  \proofstepwidestar[1,2]{T\colon\powerset(B)\to\powerset(B)}{\rIE{\FormulaRefAuto{a=b\vdash b=a}[thm]{\FormulaRefAuto{WordRecursionStageOperatorDef}[def]},\FormulaRefAuto{TypedTermGraphFunction}[thm]{10}}}
  \proofstepwidestar[1,2]{\forall H\in\powerset(B)\;T(H)=\tau(H)}{\rIE{\FormulaRefAuto{a=b\vdash b=a}[thm]{\FormulaRefAuto{WordRecursionStageOperatorDef}[def]},\FormulaRefAuto{TypedTermGraphValues}[thm]{10}}}
  \proofstepwidestar[1,2]{T\colon\powerset(B)\to\powerset(B)\land\forall H\in\powerset(B)\;T(H)=\tau(H)}{\rAI{11,12}}
\end{tabproofwide}


\Needspace{14\baselineskip}
\FormulaThmDeltaK[Zulässigkeit eines Fortsetzungspaars]{%
  \begin{aligned}[t]
    &\mathsf{WortStr}_A(W,e,s)\dsep r\colon X\times A\to X\dsep
      u\in W\dsep x\in X\dsep a\in A\vdash{}\\[-2pt]
    &(s(u,a),r(x,a))\in B
  \end{aligned}
}{WordRecursionStagePairTyping}{%
  \DeltaRow{Mengen}{A\dsep W\dsep e\dsep X\dsep u\dsep x\dsep a}\DeltaRow{Funktionen}{s\dsep r}}
\begin{tabproofwide}
  \proofstepwidestar[1]{\mathsf{WortStr}_A(W,e,s)}{\rA}
  \proofstepwidestar[2]{r\colon X\times A\to X}{\rA}
  \proofstepwidestar[3]{u\in W}{\rA}
  \proofstepwidestar[4]{x\in X}{\rA}
  \proofstepwidestar[5]{a\in A}{\rA}
  \proofstepwidestar[1,3,5]{s(u,a)\in W}{\FormulaRefAuto{WordStructureStepTyping}[thm]{1,3,5}}
  \proofstepwidestar[4,5]{(x,a)\in X\times A}{\FormulaRefAuto{(a,b)\in A\times B\eqvdash a\in A\land b\in B}[thm]{\rAI{4,5}}}
  \proofstepwidestar[2,4,5]{r(x,a)\in X}{\FormulaRefAuto{F\colon A\to B\dsep x\in A\vdash F(x)\in B}[thm]{2,7}}
  \proofstepwidestar[1,2,3,4,5]{(s(u,a),r(x,a))\in B}{\rIE{\FormulaRefAuto{a=b\vdash b=a}[thm]{\FormulaRefAuto{WordRecursionAmbientDef}[def]},\FormulaRefAuto{(a,b)\in A\times B\eqvdash a\in A\land b\in B}[thm]{\rAI{6,8}}}}
\end{tabproofwide}

\Needspace{14\baselineskip}
\FormulaThmDeltaK[Ein Paar entsteht am Anfang oder durch einen Schritt]{%
  \begin{aligned}[t]
    &\mathsf{WortStr}_A(W,e,s)\dsep x_0\in X\dsep r\colon X\times A\to X\dsep H\subseteq B\vdash{}\\[-2pt]
    &p\in T(H)\leftrightarrow
      \bigl(p=(e,x_0)\lor\exists u\in W\,\exists x\in X\,\exists a\in A\;\\[-2pt]
    &\hspace{40mm}((u,x)\in H\land p=(s(u,a),r(x,a)))\bigr)
  \end{aligned}
}{WordRecursionStageMembership}{%
  \DeltaRow{Mengen}{A\dsep W\dsep e\dsep X\dsep x_0\dsep H\dsep p\dsep u\dsep x\dsep a}
  \DeltaRow{Funktionen}{s\dsep r}
  \DeltaRow{Abkürzung}{\begin{aligned}[t]
    E(H,p)\coloneqq{}&\exists u\in W\,\exists x\in X\,\exists a\in A\;\\[-2pt]
    &((u,x)\in H\land p=(s(u,a),r(x,a)))
  \end{aligned}}}
\begin{tabproofwide}
  \proofstepwidestar[1]{\mathsf{WortStr}_A(W,e,s)}{\rA}
  \proofstepwidestar[2]{x_0\in X}{\rA}
  \proofstepwidestar[3]{r\colon X\times A\to X}{\rA}
  \proofstepwidestar[4]{H\subseteq B}{\rA}
  \proofstepwidestar[4]{H\in\powerset(B)}{\FormulaRefAuto{x\in\powerset(A)\eqvdash x\subseteq A}[thm]{4}}
  \proofstepwidestar[1,2,4]{T(H)=\tau(H)}{\rRE{5,\rUE{\rAEb{\FormulaRefAuto{WordRecursionStageOperatorFunction}[thm]{1,2}}}}}
  \proofstepwidestar[1,2,4]{p\in T(H)\leftrightarrow\bigl(p=(e,x_0)\lor(p\in B\land E(H,p))\bigr)}{\rIE{\FormulaRefAuto{a=b\vdash b=a}[thm]{6},\rIE{\FormulaRefAuto{a=b\vdash b=a}[thm]{\FormulaRefAuto{WordRecursionStageTermDef}[def]},\rLRS{\FormulaRefAuto{x\in\{a\}\eqvdash x=a}[thm],\rLRS{\FormulaRefAuto{x\in\{u\in A\mid P(u)\}\eqvdash x\in A\land P(x)}[thm],\FormulaRefAuto{x\in A\cup B\eqvdash x\in A\lor x\in B}[thm]}}}}}
  \proofstepwidestar[8]{E(H,p)}{\rA}
  \proofstepwidestar[9]{u\in W\land\exists x\in X\,\exists a\in A\;((u,x)\in H\land p=(s(u,a),r(x,a)))}{\rA}
  \proofstepwidestar[10]{x\in X\land\exists a\in A\;((u,x)\in H\land p=(s(u,a),r(x,a)))}{\rA}
  \proofstepwidestar[11]{a\in A\land((u,x)\in H\land p=(s(u,a),r(x,a)))}{\rA}
  \proofstepwidestar[1,3,9,10,11]{(s(u,a),r(x,a))\in B}{\FormulaRefAuto{WordRecursionStagePairTyping}[thm]{1,3,\rAEa{9},\rAEa{10},\rAEa{11}}}
  \proofstepwidestar[1,3,9,10,11]{p\in B}{\rIE{\FormulaRefAuto{a=b\vdash b=a}[thm]{\rAEb{\rAEb{11}}},12}}
  \proofstepwidestar[1,3,9,10]{p\in B}{\rEE{\rAEb{10},11,13}}
  \proofstepwidestar[1,3,9]{p\in B}{\rEE{\rAEb{9},10,14}}
  \proofstepwidestar[1,3,8]{p\in B}{\rEE{8,9,15}}
  \proofstepwidestar[1,3]{E(H,p)\rightarrow(p\in B\land E(H,p))}{\rRI{8,\rAI{16,8}}}
  \proofstepwidestar[18]{p\in B\land E(H,p)}{\rA}
  \proofstepwidestar[]{(p\in B\land E(H,p))\rightarrow E(H,p)}{\rRI{18,\rAEb{18}}}
  \proofstepwidestar[1,3]{(p\in B\land E(H,p))\leftrightarrow E(H,p)}{\rLRI{19,17}}
  \proofstepwidestar[1,2,3,4]{p\in T(H)\leftrightarrow\bigl(p=(e,x_0)\lor E(H,p)\bigr)}{\rLRS{20,7}}
\end{tabproofwide}

\Needspace{14\baselineskip}
\FormulaThmDeltaK[Anfang und Fortsetzung gehören zur nächsten Phase]{%
  \begin{aligned}[t]
    &\mathsf{WortStr}_A(W,e,s)\dsep x_0\in X\dsep r\colon X\times A\to X\dsep H\subseteq B\vdash{}\\[-2pt]
    &(e,x_0)\in T(H)\land
      \forall u\in W\,\forall x\in X\,\forall a\in A\;\\[-2pt]
    &\qquad((u,x)\in H\rightarrow(s(u,a),r(x,a))\in T(H))
  \end{aligned}
}{WordRecursionStageBaseAndStep}{%
  \DeltaRow{Mengen}{A\dsep W\dsep e\dsep X\dsep x_0\dsep H\dsep u\dsep x\dsep a}\DeltaRow{Funktionen}{s\dsep r}}
\begin{tabproofwide}
  \proofstepwidestar[1]{\mathsf{WortStr}_A(W,e,s)}{\rA}
  \proofstepwidestar[2]{x_0\in X}{\rA}
  \proofstepwidestar[3]{r\colon X\times A\to X}{\rA}
  \proofstepwidestar[4]{H\subseteq B}{\rA}
  \proofstepwidestar[1,2,3,4]{(e,x_0)\in T(H)}{\rRE{\rLREb{\FormulaRefAuto{WordRecursionStageMembership}[thm]{1,2,3,4}},\rOIa{\rII}}}
  \proofstepwidestar[6]{u\in W}{\rA}
  \proofstepwidestar[7]{x\in X}{\rA}
  \proofstepwidestar[8]{a\in A}{\rA}
  \proofstepwidestar[9]{(u,x)\in H}{\rA}
  \proofstepwidestar[6,7,8,9]{\exists v\in W\,\exists y\in X\,\exists b\in A\;((v,y)\in H\land(s(u,a),r(x,a))=(s(v,b),r(y,b)))}{\rEI{\rAI{6,\rEI{\rAI{7,\rEI{\rAI{8,\rAI{9,\rII}}}}}}}}
  \proofstepwidestar[1,2,3,4,6,7,8,9]{(s(u,a),r(x,a))\in T(H)}{\rRE{\rLREb{\FormulaRefAuto{WordRecursionStageMembership}[thm]{1,2,3,4}},\rOIb{10}}}
  \proofstepwidestar[1,2,3,4]{\forall u\in W\,\forall x\in X\,\forall a\in A\;((u,x)\in H\rightarrow(s(u,a),r(x,a))\in T(H))}{\rUI{\rRI{6,\rUI{\rRI{7,\rUI{\rRI{8,\rRI{9,11}}}}}}}}
  \proofstepwidestar[1,2,3,4]{\begin{aligned}[t]&(e,x_0)\in T(H)\land{}\\[-2pt]&\forall u\in W\,\forall x\in X\,\forall a\in A\;\\[-2pt]&\quad((u,x)\in H\rightarrow(s(u,a),r(x,a))\in T(H))\end{aligned}}{\rAI{5,12}}
\end{tabproofwide}

\subsubsection{Die natürlichen Zahlen zählen die Stufen}

Nun erfolgt genau die Übersetzung in die Zahlenrekursion aus Band~10:
Die dortige Wertemenge \(M\) ist hier \(\powerset(B)\), der Anfangswert
\(m_0\) ist \(\{(e,x_0)\}\), und die Schrittfunktion \(f\) ist
\(T\). Der Dedekindsche Rekursionssatz
\FormulaRefAuto{DedekindRecursionTheorem}[thm] liefert deshalb eine
Folge von Zuordnungen. Wir nennen ihre Werte \(H_n\).

\Needspace{14\baselineskip}
\FormulaDefDeltaK[Die Folge der Auswertungsstufen]{%
  \begin{aligned}[t]
    &\mathsf{WortStr}_A(W,e,s)\dsep x_0\in X\vdash{}\\[-2pt]
    &h\coloneqq\RecFun{\{(e,x_0)\}}{T}
  \end{aligned}
}{WordRecursionStagesDef}{%
  \DeltaRow{Mengen}{A\dsep W\dsep e\dsep X\dsep x_0\dsep n}
  \DeltaRow{Funktionen}{s\dsep r}\DeltaRow{Neues Symbol}{h}}

Für \(n\in\mathbb N\) verwenden wir im Folgenden die Abkürzung
\(H_n\coloneqq h(n)\).

\Needspace{14\baselineskip}
\FormulaThmDeltaK[Typisierung und Gleichungen der Stufenfolge]{%
  \begin{aligned}[t]
    &\mathsf{WortStr}_A(W,e,s)\dsep x_0\in X\vdash{}\\[-2pt]
    &h\colon\mathbb N\to\powerset(B)\land H_0=\{(e,x_0)\}\land{}\\[-2pt]
    &\forall n\in\mathbb N\;(H_n\subseteq B\land H_{\Succ(n)}=T(H_n))
  \end{aligned}
}{WordRecursionStagesEquations}{%
  \DeltaRow{Mengen}{A\dsep W\dsep e\dsep X\dsep x_0\dsep n}\DeltaRow{Funktionen}{s\dsep r}}
\begin{tabproofwide}
  \proofstepwidestar[1]{\mathsf{WortStr}_A(W,e,s)}{\rA}
  \proofstepwidestar[2]{x_0\in X}{\rA}
  \proofstepwidestar[1]{e\in W}{\rAEa{\FormulaRefAuto{WordStructureTypingAxiom}[ax]{1}}}
  \proofstepwidestar[1,2]{(e,x_0)\in B}{\rIE{\FormulaRefAuto{a=b\vdash b=a}[thm]{\FormulaRefAuto{WordRecursionAmbientDef}[def]},\FormulaRefAuto{(a,b)\in A\times B\eqvdash a\in A\land b\in B}[thm]{\rAI{3,2}}}}
  \proofstepwidestar[1,2]{\{(e,x_0)\}\in\powerset(B)}{\FormulaRefAuto{x\in\powerset(A)\eqvdash x\subseteq A}[thm]{\FormulaRefAuto{a\in A\vdash\{a\}\subseteq A}[thm]{4}}}
  \proofstepwidestar[1,2]{T\colon\powerset(B)\to\powerset(B)}{\rAEa{\FormulaRefAuto{WordRecursionStageOperatorFunction}[thm]{1,2}}}
  \proofstepwidestar[1,2]{h\colon\mathbb N\to\powerset(B)}{\rIE{\FormulaRefAuto{a=b\vdash b=a}[thm]{\FormulaRefAuto{WordRecursionStagesDef}[def]{1,2}},\FormulaRefAuto{m_0\in M\dsep f\colon M\to M\vdash\RecFun{m_0}{f}\colon\mathbb{N}\to M}[thm]{5,6}}}
  \proofstepwidestar[1,2]{H_0=\{(e,x_0)\}}{\rIE{\FormulaRefAuto{a=b\vdash b=a}[thm]{\FormulaRefAuto{WordRecursionStagesDef}[def]{1,2}},\FormulaRefAuto{m_0\in M\dsep f\colon M\to M\vdash\RecFun{m_0}{f}(0)=m_0}[thm]{5,6}}}
  \proofstepwidestar[9]{n\in\mathbb N}{\rA}
  \proofstepwidestar[1,2,9]{H_n\in\powerset(B)}{\FormulaRefAuto{F\colon A\to B\dsep x\in A\vdash F(x)\in B}[thm]{7,9}}
  \proofstepwidestar[1,2,9]{H_n\subseteq B}{\FormulaRefAuto{x\in\powerset(A)\eqvdash x\subseteq A}[thm]{10}}
  \proofstepwidestar[1,2,9]{H_{\Succ(n)}=T(H_n)}{\rIE{\FormulaRefAuto{a=b\vdash b=a}[thm]{\FormulaRefAuto{WordRecursionStagesDef}[def]{1,2}},\FormulaRefAuto{m_0\in M\dsep f\colon M\to M\dsep n\in\mathbb{N}\vdash\RecFun{m_0}{f}(\Succ(n))=f(\RecFun{m_0}{f}(n))}[thm]{5,6,9}}}
  \proofstepwidestar[1,2]{\forall n\in\mathbb N\;(H_n\subseteq B\land H_{\Succ(n)}=T(H_n))}{\rUI{\rRI{9,\rAI{11,12}}}}
  \proofstepwidestar[1,2]{\begin{aligned}[t]&h\colon\mathbb N\to\powerset(B)\land H_0=\{(e,x_0)\}\land{}\\[-2pt]&\forall n\in\mathbb N\;(H_n\subseteq B\land H_{\Succ(n)}=T(H_n))\end{aligned}}{\rAI{7,\rAI{8,13}}}
\end{tabproofwide}

\(H_0\) ordnet dem Anfangswort seinen Wert zu. \(H_1\) enthält
zusätzlich die Werte der Einbuchstabenwörter, \(H_2\) die der nächsten
Fortsetzungen. Diese Beschreibung dient der Anschauung; im Beweis
arbeiten wir allein mit den beiden Stufengleichungen. Bei unendlichem
Alphabet kann bereits \(H_1\) unendlich sein. Bis zu jeder festen
Stufe sind nur endlich viele Aufbauphasen nötig; die Zahl ihrer
Paare muss dagegen nicht endlich sein.

\subsubsection{Spätere Stufen bewahren frühere Zuordnungen}

Obwohl \(T(H)\) nur das Anfangspaar und die Fortsetzungen enthält,
gehen in der konstruierten Folge keine Paare verloren. Der Grund ist
derselbe Induktionsgedanke wie bei wachsenden Folgen natürlicher
Anfangsabschnitte: Die erste Stufe ist in der zweiten enthalten, und
ein Einschluss bleibt unter einem weiteren Aufbau erhalten.

\Needspace{14\baselineskip}
\FormulaThmDeltaK[Die Aufbauphase erhält Einschlüsse]{%
  \begin{aligned}[t]
    &\mathsf{WortStr}_A(W,e,s)\dsep x_0\in X\dsep r\colon X\times A\to X\dsep{}\\[-2pt]
    &H\subseteq B\dsep K\subseteq B\dsep H\subseteq K\vdash T(H)\subseteq T(K)
  \end{aligned}
}{WordRecursionStageMonotone}{%
  \DeltaRow{Mengen}{A\dsep W\dsep e\dsep X\dsep x_0\dsep H\dsep K\dsep p\dsep u\dsep x\dsep a}\DeltaRow{Funktionen}{s\dsep r}}
\begin{tabproofwide}
  \proofstepwidestar[1]{\mathsf{WortStr}_A(W,e,s)}{\rA}
  \proofstepwidestar[2]{x_0\in X}{\rA}
  \proofstepwidestar[3]{r\colon X\times A\to X}{\rA}
  \proofstepwidestar[4]{H\subseteq B}{\rA}
  \proofstepwidestar[5]{K\subseteq B}{\rA}
  \proofstepwidestar[6]{H\subseteq K}{\rA}
  \proofstepwidestar[7]{p\in T(H)}{\rA}
  \proofstepwidestar[1,2,3,4,7]{p=(e,x_0)\lor\exists u\in W\,\exists x\in X\,\exists a\in A\;((u,x)\in H\land p=(s(u,a),r(x,a)))}{\rRE{7,\rLREa{\FormulaRefAuto{WordRecursionStageMembership}[thm]{1,2,3,4}}}}
  \proofstepwidestar[9]{p=(e,x_0)}{\rA}
  \proofstepwidestar[1,2,3,5,9]{p\in T(K)}{\rIE{\FormulaRefAuto{a=b\vdash b=a}[thm]{9},\rAEa{\FormulaRefAuto{WordRecursionStageBaseAndStep}[thm]{1,2,3,5}}}}
  \proofstepwidestar[11]{\exists u\in W\,\exists x\in X\,\exists a\in A\;((u,x)\in H\land p=(s(u,a),r(x,a)))}{\rA}
  \proofstepwidestar[12]{u\in W\land\exists x\in X\,\exists a\in A\;((u,x)\in H\land p=(s(u,a),r(x,a)))}{\rA}
  \proofstepwidestar[13]{x\in X\land\exists a\in A\;((u,x)\in H\land p=(s(u,a),r(x,a)))}{\rA}
  \proofstepwidestar[14]{a\in A\land((u,x)\in H\land p=(s(u,a),r(x,a)))}{\rA}
  \proofstepwidestar[6,14]{(u,x)\in K}{\FormulaRefAuto{A\subseteq B\dsep x\in A\vdash x\in B}[thm]{6,\rAEa{\rAEb{14}}}}
  \proofstepwidestar[6,12,13,14]{\exists u\in W\,\exists x\in X\,\exists a\in A\;((u,x)\in K\land p=(s(u,a),r(x,a)))}{\rEI{\rAI{\rAEa{12},\rEI{\rAI{\rAEa{13},\rEI{\rAI{\rAEa{14},\rAI{15,\rAEb{\rAEb{14}}}}}}}}}}
  \proofstepwidestar[1,2,3,5,6,12,13,14]{p\in T(K)}{\rRE{\rLREb{\FormulaRefAuto{WordRecursionStageMembership}[thm]{1,2,3,5}},\rOIb{16}}}
  \proofstepwidestar[1,2,3,5,6,12,13]{p\in T(K)}{\rEE{\rAEb{13},14,17}}
  \proofstepwidestar[1,2,3,5,6,12]{p\in T(K)}{\rEE{\rAEb{12},13,18}}
  \proofstepwidestar[1,2,3,5,6,11]{p\in T(K)}{\rEE{11,12,19}}
  \proofstepwidestar[1,2,3,4,5,6,7]{p\in T(K)}{\rOE{8,9,10,11,20}}
  \proofstepwidestar[1,2,3,4,5,6]{T(H)\subseteq T(K)}{\FormulaRefAuto{SubsetIntroductionRule}[thm]{[7]\vdots21}}
\end{tabproofwide}

\Needspace{14\baselineskip}
\FormulaThmDeltaK[Jede Stufe liegt in der nächsten]{%
  \begin{aligned}[t]
    &\mathsf{WortStr}_A(W,e,s)\dsep x_0\in X\dsep r\colon X\times A\to X\vdash{}\\[-2pt]
    &\forall n\in\mathbb N\;H_n\subseteq H_{\Succ(n)}
  \end{aligned}
}{WordRecursionStagesIncreasing}{%
  \DeltaRow{Mengen}{A\dsep W\dsep e\dsep X\dsep x_0\dsep n\dsep k}
  \DeltaRow{Funktionen}{s\dsep r}\DeltaRow{Abkürzung}{P(n)\coloneqq H_n\subseteq H_{\Succ(n)}}}
\begin{tabproofwide}
  \proofstepwidestar[1]{\mathsf{WortStr}_A(W,e,s)}{\rA}
  \proofstepwidestar[2]{x_0\in X}{\rA}
  \proofstepwidestar[3]{r\colon X\times A\to X}{\rA}
  \proofstepwidestar[1,2]{H_0=\{(e,x_0)\}\land\forall n\in\mathbb N\;(H_n\subseteq B\land H_{\Succ(n)}=T(H_n))}{\rAEb{\FormulaRefAuto{WordRecursionStagesEquations}[thm]{1,2}}}
  \proofstepwidestar[]{0\in\mathbb N}{\FormulaRefAuto{PeanoZero}[ax]}
  \proofstepwidestar[1,2]{H_0\subseteq B\land H_{\Succ(0)}=T(H_0)}{\rRE{5,\rUE{\rAEb{4}}}}
  \proofstepwidestar[1,2,3]{(e,x_0)\in T(H_0)}{\rAEa{\FormulaRefAuto{WordRecursionStageBaseAndStep}[thm]{1,2,3,\rAEa{6}}}}
  \proofstepwidestar[1,2,3]{P(0)}{\rIE{\FormulaRefAuto{a=b\vdash b=a}[thm]{\rAEa{4}},\FormulaRefAuto{a=b\vdash b=a}[thm]{\rAEb{6}},\FormulaRefAuto{a\in A\vdash\{a\}\subseteq A}[thm]{7}}}
  \proofstepwidestar[9]{n\in\mathbb N}{\rA}
  \proofstepwidestar[10]{P(n)}{\rA}
  \proofstepwidestar[9]{\Succ(n)\in\mathbb N}{\FormulaRefAuto{PeanoSuccClosure}[ax]{9}}
  \proofstepwidestar[1,2,9]{H_n\subseteq B\land H_{\Succ(n)}=T(H_n)}{\rRE{9,\rUE{\rAEb{4}}}}
  \proofstepwidestar[1,2,9]{H_{\Succ(n)}\subseteq B\land H_{\Succ(\Succ(n))}=T(H_{\Succ(n)})}{\rRE{11,\rUE{\rAEb{4}}}}
  \proofstepwidestar[1,2,3,9,10]{T(H_n)\subseteq T(H_{\Succ(n)})}{\FormulaRefAuto{WordRecursionStageMonotone}[thm]{1,2,3,\rAEa{12},\rAEa{13},10}}
  \proofstepwidestar[1,2,3,9,10]{P(\Succ(n))}{\rIE{\FormulaRefAuto{a=b\vdash b=a}[thm]{\rAEb{12}},\FormulaRefAuto{a=b\vdash b=a}[thm]{\rAEb{13}},14}}
  \proofstepwidestar[1,2,3]{\forall n\in\mathbb N\;(P(n)\rightarrow P(\Succ(n)))}{\rUI{\rRI{9,\rRI{10,15}}}}
  \proofstepwidestar[17]{k\in\mathbb N}{\rA}
  \proofstepwidestar[1,2,3,17]{P(k)}{\FormulaRefAuto{PeanoInductionRule}[ax]{8,[9,10]\vdots15,17}}
  \proofstepwidestar[1,2,3]{\forall n\in\mathbb N\;H_n\subseteq H_{\Succ(n)}}{\rUI{\rRI{17,18}}}
\end{tabproofwide}

Der Monotoniesatz aus Band~10 ist in der Schreibweise \(n+1\)
formuliert. Mit \(n+1=\Succ(n)\) übersetzen wir zunächst den
soeben bewiesenen Stufeneinschluss in diese Schreibweise. Die
Aufbauphase selbst bleibt dabei dieselbe.

\Needspace{14\baselineskip}
\FormulaThmDeltaK[Frühere Stufen liegen in allen späteren]{%
  \begin{aligned}[t]
    &\mathsf{WortStr}_A(W,e,s)\dsep x_0\in X\dsep r\colon X\times A\to X\dsep{}\\[-2pt]
    &n\in\mathbb N\dsep m\in\mathbb N\dsep n\leq m\vdash H_n\subseteq H_m
  \end{aligned}
}{WordRecursionStagesChain}{%
  \DeltaRow{Mengen}{A\dsep W\dsep e\dsep X\dsep x_0\dsep n\dsep m\dsep k}\DeltaRow{Funktionen}{s\dsep r}}
\begin{tabproofwide}
  \proofstepwidestar[1]{\mathsf{WortStr}_A(W,e,s)}{\rA}
  \proofstepwidestar[2]{x_0\in X}{\rA}
  \proofstepwidestar[3]{r\colon X\times A\to X}{\rA}
  \proofstepwidestar[4]{n\in\mathbb N}{\rA}
  \proofstepwidestar[5]{m\in\mathbb N}{\rA}
  \proofstepwidestar[6]{n\leq m}{\rA}
  \proofstepwidestar[1,2,3]{\forall k\in\mathbb N\;H_k\subseteq H_{\Succ(k)}}{\FormulaRefAuto{WordRecursionStagesIncreasing}[thm]{1,2,3}}
  \proofstepwidestar[8]{k\in\mathbb N}{\rA}
  \proofstepwidestar[1,2,3,8]{H_k\subseteq H_{\Succ(k)}}{\rRE{8,\rUE{7}}}
  \proofstepwidestar[8]{k+1=\Succ(k)}{\FormulaRefAuto{PeanoAddOneSucc}[thm]{8}}
  \proofstepwidestar[1,2,3,8]{H_k\subseteq H_{k+1}}{\rIE{\FormulaRefAuto{a=b\vdash b=a}[thm]{10},9}}
  \proofstepwidestar[1,2,3]{\forall k\in\mathbb N\;H_k\subseteq H_{k+1}}{\rUI{\rRI{8,11}}}
  \proofstepwidestar[4,5,6]{n<m\lor n=m}{\FormulaRefAuto{m\in\mathbb{N}\dsep n\in\mathbb{N}\dsep m\leq n\vdash m<n\lor m=n}[thm]{4,5,6}}
  \proofstepwidestar[14]{n<m}{\rA}
  \proofstepwidestar[1,2,3,4]{H_n\subseteq H_{n+1}}{\rRE{4,\rUE{12}}}
  \proofstepwidestar[1,2,3,4,5,14]{H_{n+1}\subseteq H_m}{\FormulaRefAuto{B04PeanoStrictMonotoneAddOne}[thm]{12,4,5,14}}
  \proofstepwidestar[1,2,3,4,5,14]{H_n\subseteq H_m}{\FormulaRefAuto{A\subseteq B,B\subseteq C\vdash A\subseteq C}[thm]{15,16}}
  \proofstepwidestar[18]{n=m}{\rA}
  \proofstepwidestar[18]{H_n\subseteq H_m}{\rIE{18,\FormulaRefAuto{A\subseteq A}[thm]}}
  \proofstepwidestar[1,2,3,4,5,6]{H_n\subseteq H_m}{\rOE{13,14,17,18,19}}
\end{tabproofwide}

\subsubsection{Jedes erfasste Wort erhält nur einen Wert}

Nun treten die Unterschiede zwischen Zahlen und Wörtern im Beweis
ausdrücklich hervor. Bei Zahlen verhindert das Nullaxiom eine Kollision
zwischen Anfang und Nachfolger; hier leistet W1 dasselbe für das
Anfangswort. Die Injektivität des Zahlennachfolgers gewinnt den Vorgänger
zurück. W2 gewinnt hier sowohl den Vorgänger als auch den angefügten
Buchstaben zurück. Deshalb können verschiedene Berechnungswege einem
Wort keinen widersprüchlichen Wert geben.

Wir kürzen lediglich die gewöhnliche Rechtseindeutigkeit einer Relation
ab. Sie fordert noch nicht, dass jedes Wort einen Wert besitzt.

\Needspace{14\baselineskip}
\FormulaDefDeltaK[Ein Wert je erfasstem Argument]{%
  R(H)\coloneqq\forall t\,\forall y\,\forall z\;
    (((t,y)\in H\land(t,z)\in H)\rightarrow y=z)
}{WordRecursionRightUniqueDef}{%
  \DeltaRow{Mengen}{H\dsep t\dsep y\dsep z}\DeltaRow{Neues Symbol}{R(H)}}

\Needspace{14\baselineskip}
\FormulaThmDeltaK[Die Anfangszuordnung ist rechtseindeutig]{%
  R(\{(e,x_0)\})
}{WordRecursionSingletonRightUnique}{%
  \DeltaRow{Mengen}{e\dsep x_0\dsep t\dsep y\dsep z}}
\begin{tabproofwide}
  \proofstepwidestar[1]{(t,y)\in\{(e,x_0)\}\land(t,z)\in\{(e,x_0)\}}{\rA}
  \proofstepwidestar[1]{(t,y)=(e,x_0)}{\FormulaRefAuto{x\in\{a\}\eqvdash x=a}[thm]{\rAEa{1}}}
  \proofstepwidestar[1]{(t,z)=(e,x_0)}{\FormulaRefAuto{x\in\{a\}\eqvdash x=a}[thm]{\rAEb{1}}}
  \proofstepwidestar[1]{y=x_0}{\rAEb{\FormulaRefAuto{(a,b)=(c,d)\eqvdash a=c\land b=d}[thm]{2}}}
  \proofstepwidestar[1]{z=x_0}{\rAEb{\FormulaRefAuto{(a,b)=(c,d)\eqvdash a=c\land b=d}[thm]{3}}}
  \proofstepwidestar[1]{y=z}{\FormulaRefAuto{a=b\dsep c=b\vdash a=c}[thm]{4,5}}
  \proofstepwidestar[]{R(\{(e,x_0)\})}{\FormulaRefAuto{WordRecursionRightUniqueDef}[def]{\rUI{\rUI{\rUI{\rRI{1,6}}}}}}
\end{tabproofwide}

\Needspace{14\baselineskip}
\FormulaThmDeltaK[Ein Schritt kann den Anfangswert nicht verändern]{%
  \begin{aligned}[t]
    &\mathsf{WortStr}_A(W,e,s)\dsep x_0\in X\dsep r\colon X\times A\to X\dsep{}\\[-2pt]
    &H\subseteq B\dsep(e,z)\in T(H)\vdash z=x_0
  \end{aligned}
}{WordRecursionStageZeroValue}{%
  \DeltaRow{Mengen}{A\dsep W\dsep e\dsep X\dsep x_0\dsep H\dsep z\dsep u\dsep x\dsep a}\DeltaRow{Funktionen}{s\dsep r}}
\begin{tabproofwide}
  \proofstepwidestar[1]{\mathsf{WortStr}_A(W,e,s)}{\rA}
  \proofstepwidestar[2]{x_0\in X}{\rA}
  \proofstepwidestar[3]{r\colon X\times A\to X}{\rA}
  \proofstepwidestar[4]{H\subseteq B}{\rA}
  \proofstepwidestar[5]{(e,z)\in T(H)}{\rA}
  \proofstepwidestar[1,2,3,4,5]{(e,z)=(e,x_0)\lor\exists u\in W\,\exists x\in X\,\exists a\in A\;((u,x)\in H\land(e,z)=(s(u,a),r(x,a)))}{\rRE{5,\rLREa{\FormulaRefAuto{WordRecursionStageMembership}[thm]{1,2,3,4}}}}
  \proofstepwidestar[7]{(e,z)=(e,x_0)}{\rA}
  \proofstepwidestar[7]{z=x_0}{\rAEb{\FormulaRefAuto{(a,b)=(c,d)\eqvdash a=c\land b=d}[thm]{7}}}
  \proofstepwidestar[9]{\exists u\in W\,\exists x\in X\,\exists a\in A\;((u,x)\in H\land(e,z)=(s(u,a),r(x,a)))}{\rA}
  \proofstepwidestar[10]{u\in W\land\exists x\in X\,\exists a\in A\;((u,x)\in H\land(e,z)=(s(u,a),r(x,a)))}{\rA}
  \proofstepwidestar[11]{x\in X\land\exists a\in A\;((u,x)\in H\land(e,z)=(s(u,a),r(x,a)))}{\rA}
  \proofstepwidestar[12]{a\in A\land((u,x)\in H\land(e,z)=(s(u,a),r(x,a)))}{\rA}
  \proofstepwidestar[12]{e=s(u,a)}{\rAEa{\FormulaRefAuto{(a,b)=(c,d)\eqvdash a=c\land b=d}[thm]{\rAEb{\rAEb{12}}}}}
  \proofstepwidestar[1,10,12]{s(u,a)\ne e}{\FormulaRefAuto{WordStructureZeroAxiom}[ax]{1,\rAEa{10},\rAEa{12}}}
  \proofstepwidestar[1,10,12]{\bot}{\rBI{\FormulaRefAuto{a=b\vdash b=a}[thm]{13},14}}
  \proofstepwidestar[1,10,12]{z=x_0}{\FormulaRefAuto{\bot\vdash P}[thm]{15}}
  \proofstepwidestar[1,10,11]{z=x_0}{\rEE{\rAEb{11},12,16}}
  \proofstepwidestar[1,10]{z=x_0}{\rEE{\rAEb{10},11,17}}
  \proofstepwidestar[1,9]{z=x_0}{\rEE{9,10,18}}
  \proofstepwidestar[1,2,3,4,5]{z=x_0}{\rOE{6,7,8,9,19}}
\end{tabproofwide}

\Needspace{14\baselineskip}
\FormulaThmDeltaK[Ein Nachfolgerwert wird vom Vorgänger eindeutig festgelegt]{%
  \begin{aligned}[t]
    &\mathsf{WortStr}_A(W,e,s)\dsep x_0\in X\dsep r\colon X\times A\to X\dsep{}\\[-2pt]
    &H\subseteq B\dsep R(H)\dsep u\in W\dsep x\in X\dsep a\in A\dsep{}\\[-2pt]
    &(u,x)\in H\dsep(s(u,a),z)\in T(H)\vdash z=r(x,a)
  \end{aligned}
}{WordRecursionStageSuccessorValue}{%
  \DeltaRow{Mengen}{A\dsep W\dsep e\dsep X\dsep x_0\dsep H\dsep u\dsep x\dsep a\dsep z\dsep v\dsep y\dsep b}
  \DeltaRow{Funktionen}{s\dsep r}}
\begin{tabproofwide}
  \proofstepwidestar[1]{\mathsf{WortStr}_A(W,e,s)}{\rA}
  \proofstepwidestar[2]{x_0\in X}{\rA}
  \proofstepwidestar[3]{r\colon X\times A\to X}{\rA}
  \proofstepwidestar[4]{H\subseteq B}{\rA}
  \proofstepwidestar[5]{R(H)}{\rA}
  \proofstepwidestar[6]{u\in W}{\rA}
  \proofstepwidestar[7]{x\in X}{\rA}
  \proofstepwidestar[8]{a\in A}{\rA}
  \proofstepwidestar[9]{(u,x)\in H}{\rA}
  \proofstepwidestar[10]{(s(u,a),z)\in T(H)}{\rA}
  \proofstepwidestar[1,2,3,4,10]{\begin{aligned}[t]&(s(u,a),z)=(e,x_0)\lor{}\\[-2pt]&\exists v\in W\,\exists y\in X\,\exists b\in A\;\\[-2pt]&\quad((v,y)\in H\land(s(u,a),z)=(s(v,b),r(y,b)))\end{aligned}}{\rRE{10,\rLREa{\FormulaRefAuto{WordRecursionStageMembership}[thm]{1,2,3,4}}}}
  \proofstepwidestar[12]{(s(u,a),z)=(e,x_0)}{\rA}
  \proofstepwidestar[12]{s(u,a)=e}{\rAEa{\FormulaRefAuto{(a,b)=(c,d)\eqvdash a=c\land b=d}[thm]{12}}}
  \proofstepwidestar[1,6,8]{s(u,a)\ne e}{\FormulaRefAuto{WordStructureZeroAxiom}[ax]{1,6,8}}
  \proofstepwidestar[1,6,8,12]{z=r(x,a)}{\FormulaRefAuto{\bot\vdash P}[thm]{\rBI{13,14}}}
  \proofstepwidestar[16]{\exists v\in W\,\exists y\in X\,\exists b\in A\;((v,y)\in H\land(s(u,a),z)=(s(v,b),r(y,b)))}{\rA}
  \proofstepwidestar[17]{v\in W\land\exists y\in X\,\exists b\in A\;((v,y)\in H\land(s(u,a),z)=(s(v,b),r(y,b)))}{\rA}
  \proofstepwidestar[18]{y\in X\land\exists b\in A\;((v,y)\in H\land(s(u,a),z)=(s(v,b),r(y,b)))}{\rA}
  \proofstepwidestar[19]{b\in A\land((v,y)\in H\land(s(u,a),z)=(s(v,b),r(y,b)))}{\rA}
  \proofstepwidestar[19]{s(u,a)=s(v,b)\land z=r(y,b)}{\FormulaRefAuto{(a,b)=(c,d)\eqvdash a=c\land b=d}[thm]{\rAEb{\rAEb{19}}}}
  \proofstepwidestar[1,6,8,17,19]{u=v\land a=b}{\FormulaRefAuto{WordStructureInjectivityAxiom}[ax]{1,6,\rAEa{17},8,\rAEa{19},\rAEa{20}}}
  \proofstepwidestar[1,6,8,17,19]{(u,y)\in H}{\rIE{\FormulaRefAuto{a=b\vdash b=a}[thm]{\rAEa{21}},\rAEa{\rAEb{19}}}}
  \proofstepwidestar[1,5,6,8,9,17,19]{x=y}{\rRE{\rAI{9,22},\rUE{\rUE{\rUE{\FormulaRefAuto{WordRecursionRightUniqueDef}[def]{5}}}}}}
  \proofstepwidestar[1,5,6,8,9,17,19]{z=r(x,a)}{\rIE{\FormulaRefAuto{a=b\vdash b=a}[thm]{23},\FormulaRefAuto{a=b\vdash b=a}[thm]{\rAEb{21}},\rAEb{20}}}
  \proofstepwidestar[1,5,6,8,9,17,18]{z=r(x,a)}{\rEE{\rAEb{18},19,24}}
  \proofstepwidestar[1,5,6,8,9,17]{z=r(x,a)}{\rEE{\rAEb{17},18,25}}
  \proofstepwidestar[1,5,6,8,9,16]{z=r(x,a)}{\rEE{16,17,26}}
  \proofstepwidestar[1,2,3,4,5,6,8,9,10]{z=r(x,a)}{\rOE{11,12,15,16,27}}
\end{tabproofwide}

\Needspace{14\baselineskip}
\FormulaThmDeltaK[Rechtseindeutigkeit bleibt in der nächsten Phase erhalten]{%
  \begin{aligned}[t]
    &\mathsf{WortStr}_A(W,e,s)\dsep x_0\in X\dsep r\colon X\times A\to X\dsep{}\\[-2pt]
    &H\subseteq B\dsep R(H)\vdash R(T(H))
  \end{aligned}
}{WordRecursionStageRightUnique}{%
  \DeltaRow{Mengen}{A\dsep W\dsep e\dsep X\dsep x_0\dsep H\dsep t\dsep z\dsep c\dsep u\dsep x\dsep a}
  \DeltaRow{Funktionen}{s\dsep r}}
\begin{tabproofwide}
  \proofstepwidestar[1]{\mathsf{WortStr}_A(W,e,s)}{\rA}
  \proofstepwidestar[2]{x_0\in X}{\rA}
  \proofstepwidestar[3]{r\colon X\times A\to X}{\rA}
  \proofstepwidestar[4]{H\subseteq B}{\rA}
  \proofstepwidestar[5]{R(H)}{\rA}
  \proofstepwidestar[6]{(t,z)\in T(H)\land(t,c)\in T(H)}{\rA}
  \proofstepwidestar[1,2,3,4,6]{(t,z)=(e,x_0)\lor\exists u\in W\,\exists x\in X\,\exists a\in A\;((u,x)\in H\land(t,z)=(s(u,a),r(x,a)))}{\rRE{\rAEa{6},\rLREa{\FormulaRefAuto{WordRecursionStageMembership}[thm]{1,2,3,4}}}}
  \proofstepwidestar[8]{(t,z)=(e,x_0)}{\rA}
  \proofstepwidestar[8]{t=e\land z=x_0}{\FormulaRefAuto{(a,b)=(c,d)\eqvdash a=c\land b=d}[thm]{8}}
  \proofstepwidestar[6,8]{(e,c)\in T(H)}{\rIE{\rAEa{9},\rAEb{6}}}
  \proofstepwidestar[1,2,3,4,6,8]{c=x_0}{\FormulaRefAuto{WordRecursionStageZeroValue}[thm]{1,2,3,4,10}}
  \proofstepwidestar[1,2,3,4,6,8]{z=c}{\FormulaRefAuto{a=b\dsep c=b\vdash a=c}[thm]{\rAEb{9},11}}
  \proofstepwidestar[13]{\exists u\in W\,\exists x\in X\,\exists a\in A\;((u,x)\in H\land(t,z)=(s(u,a),r(x,a)))}{\rA}
  \proofstepwidestar[14]{u\in W\land\exists x\in X\,\exists a\in A\;((u,x)\in H\land(t,z)=(s(u,a),r(x,a)))}{\rA}
  \proofstepwidestar[15]{x\in X\land\exists a\in A\;((u,x)\in H\land(t,z)=(s(u,a),r(x,a)))}{\rA}
  \proofstepwidestar[16]{a\in A\land((u,x)\in H\land(t,z)=(s(u,a),r(x,a)))}{\rA}
  \proofstepwidestar[16]{t=s(u,a)\land z=r(x,a)}{\FormulaRefAuto{(a,b)=(c,d)\eqvdash a=c\land b=d}[thm]{\rAEb{\rAEb{16}}}}
  \proofstepwidestar[6,16]{(s(u,a),c)\in T(H)}{\rIE{\rAEa{17},\rAEb{6}}}
  \proofstepwidestar[1,2,3,4,5,6,14,15,16]{c=r(x,a)}{\FormulaRefAuto{WordRecursionStageSuccessorValue}[thm]{1,2,3,4,5,\rAEa{14},\rAEa{15},\rAEa{16},\rAEa{\rAEb{16}},18}}
  \proofstepwidestar[1,2,3,4,5,6,14,15,16]{z=c}{\FormulaRefAuto{a=b\dsep c=b\vdash a=c}[thm]{\rAEb{17},19}}
  \proofstepwidestar[1,2,3,4,5,6,14,15]{z=c}{\rEE{\rAEb{15},16,20}}
  \proofstepwidestar[1,2,3,4,5,6,14]{z=c}{\rEE{\rAEb{14},15,21}}
  \proofstepwidestar[1,2,3,4,5,6,13]{z=c}{\rEE{13,14,22}}
  \proofstepwidestar[1,2,3,4,5,6]{z=c}{\rOE{7,8,12,13,23}}
  \proofstepwidestar[1,2,3,4,5]{R(T(H))}{\FormulaRefAuto{WordRecursionRightUniqueDef}[def]{\rUI{\rUI{\rUI{\rRI{6,24}}}}}}
\end{tabproofwide}

Damit sind Anfang und Schritt einer gewöhnlichen Zahleninduktion
bewiesen. Sie zeigt die Rechtseindeutigkeit aller Stufen. Die
Wortinduktion brauchen wir erst anschließend für die andere Frage,
ob jedes Wort überhaupt in einer Stufe vorkommt.

\Needspace{14\baselineskip}
\FormulaThmDeltaK[Alle Stufen sind rechtseindeutig]{%
  \begin{aligned}[t]
    &\mathsf{WortStr}_A(W,e,s)\dsep x_0\in X\dsep r\colon X\times A\to X\vdash{}\\[-2pt]
    &\forall n\in\mathbb N\;R(H_n)
  \end{aligned}
}{WordRecursionStagesRightUnique}{%
  \DeltaRow{Mengen}{A\dsep W\dsep e\dsep X\dsep x_0\dsep n\dsep k}\DeltaRow{Funktionen}{s\dsep r}}
\begin{tabproofwide}
  \proofstepwidestar[1]{\mathsf{WortStr}_A(W,e,s)}{\rA}
  \proofstepwidestar[2]{x_0\in X}{\rA}
  \proofstepwidestar[3]{r\colon X\times A\to X}{\rA}
  \proofstepwidestar[1,2]{H_0=\{(e,x_0)\}\land\forall n\in\mathbb N\;(H_n\subseteq B\land H_{\Succ(n)}=T(H_n))}{\rAEb{\FormulaRefAuto{WordRecursionStagesEquations}[thm]{1,2}}}
  \proofstepwidestar[1,2]{R(H_0)}{\rIE{\FormulaRefAuto{a=b\vdash b=a}[thm]{\rAEa{4}},\FormulaRefAuto{WordRecursionSingletonRightUnique}[thm]}}
  \proofstepwidestar[6]{n\in\mathbb N}{\rA}
  \proofstepwidestar[7]{R(H_n)}{\rA}
  \proofstepwidestar[1,2,6]{H_n\subseteq B\land H_{\Succ(n)}=T(H_n)}{\rRE{6,\rUE{\rAEb{4}}}}
  \proofstepwidestar[1,2,3,6,7]{R(T(H_n))}{\FormulaRefAuto{WordRecursionStageRightUnique}[thm]{1,2,3,\rAEa{8},7}}
  \proofstepwidestar[1,2,3,6,7]{R(H_{\Succ(n)})}{\rIE{\FormulaRefAuto{a=b\vdash b=a}[thm]{\rAEb{8}},9}}
  \proofstepwidestar[1,2,3]{\forall n\in\mathbb N\;(R(H_n)\rightarrow R(H_{\Succ(n)}))}{\rUI{\rRI{6,\rRI{7,10}}}}
  \proofstepwidestar[12]{k\in\mathbb N}{\rA}
  \proofstepwidestar[1,2,3,12]{R(H_k)}{\FormulaRefAuto{PeanoInductionRule}[ax]{5,[6,7]\vdots10,12}}
  \proofstepwidestar[1,2,3]{\forall n\in\mathbb N\;R(H_n)}{\rUI{\rRI{12,13}}}
\end{tabproofwide}

\subsubsection{Die Vereinigung erfasst alle Wörter}

Wir fassen jetzt alle Stufen zusammen. Ein Paar gehört zur Vereinigung,
sobald es in irgendeiner Stufe vorkommt. Zwei Stufen sind nach ihren
natürlichen Indizes vergleichbar; deshalb liegen zwei zu prüfende Paare
immer gemeinsam in einer späteren Stufe. Deren Rechtseindeutigkeit
überträgt sich auf die Vereinigung.

\Needspace{14\baselineskip}
\FormulaDefDeltaK[Vereinigung der Auswertungsstufen]{%
  \mathsf{WortStr}_A(W,e,s)\dsep x_0\in X\vdash
  G\coloneqq\bigcup h[\mathbb N]
}{WordRecursionStageUnionDef}{%
  \DeltaRow{Mengen}{A\dsep W\dsep e\dsep X\dsep x_0}\DeltaRow{Funktionen}{s\dsep r}
  \DeltaRow{Neues Symbol}{G}}

\Needspace{14\baselineskip}
\FormulaThmDeltaK[Zugehörigkeit zur Vereinigung der Stufen]{%
  \begin{aligned}[t]
    &\mathsf{WortStr}_A(W,e,s)\dsep x_0\in X\vdash{}\\[-2pt]
    &p\in G\leftrightarrow\exists n\in\mathbb N\;p\in H_n
  \end{aligned}
}{WordRecursionStageUnionMembership}{%
  \DeltaRow{Mengen}{A\dsep W\dsep e\dsep X\dsep x_0\dsep p\dsep K\dsep n}\DeltaRow{Funktionen}{s\dsep r}}
\begin{tabproofwide}
  \proofstepwidestar[1]{\mathsf{WortStr}_A(W,e,s)}{\rA}
  \proofstepwidestar[2]{x_0\in X}{\rA}
  \proofstepwidestar[1,2]{h\colon\mathbb N\to\powerset(B)}{\rAEa{\FormulaRefAuto{WordRecursionStagesEquations}[thm]{1,2}}}
  \proofstepwidestar[4]{p\in G}{\rA}
  \proofstepwidestar[1,2,4]{\exists K\in h[\mathbb N]\;p\in K}{\FormulaRefAuto{x\in\bigcup A\eqvdash\exists B\in A\;(x\in B)}[thm]{\rIE{\FormulaRefAuto{WordRecursionStageUnionDef}[def]{1,2},4}}}
  \proofstepwidestar[6]{K\in h[\mathbb N]\land p\in K}{\rA}
  \proofstepwidestar[1,2,6]{\exists n\in\mathbb N\;K=H_n}{\FormulaRefAuto{F\colon A\to B\dsep y\in F[A]\vdash\exists x\in A\,y=F(x)}[thm]{3,\rAEa{6}}}
  \proofstepwidestar[8]{n\in\mathbb N\land K=H_n}{\rA}
  \proofstepwidestar[6,8]{p\in H_n}{\rIE{\rAEb{8},\rAEb{6}}}
  \proofstepwidestar[6,8]{\exists n\in\mathbb N\;p\in H_n}{\rEI{\rAI{\rAEa{8},9}}}
  \proofstepwidestar[1,2,6]{\exists n\in\mathbb N\;p\in H_n}{\rEE{7,8,10}}
  \proofstepwidestar[1,2,4]{\exists n\in\mathbb N\;p\in H_n}{\rEE{5,6,11}}
  \proofstepwidestar[13]{\exists n\in\mathbb N\;p\in H_n}{\rA}
  \proofstepwidestar[14]{n\in\mathbb N\land p\in H_n}{\rA}
  \proofstepwidestar[14]{\exists k\in\mathbb N\;H_n=H_k}{\rEI{\rAI{\rAEa{14},\rII}}}
  \proofstepwidestar[1,2,14]{H_n\in h[\mathbb N]}{\FormulaRefAuto{F\colon A\to B\dsep\exists x\in A\,y=F(x)\vdash y\in F[A]}[thm]{3,15}}
  \proofstepwidestar[1,2,14]{p\in\bigcup h[\mathbb N]}{\FormulaRefAuto{x\in\bigcup A\eqvdash\exists B\in A\;(x\in B)}[thm]{\rEI{\rAI{16,\rAEb{14}}}}}
  \proofstepwidestar[1,2,14]{p\in G}{\rIE{\FormulaRefAuto{a=b\vdash b=a}[thm]{\FormulaRefAuto{WordRecursionStageUnionDef}[def]{1,2}},17}}
  \proofstepwidestar[1,2,13]{p\in G}{\rEE{13,14,18}}
  \proofstepwidestar[1,2]{p\in G\leftrightarrow\exists n\in\mathbb N\;p\in H_n}{\rLRI{\rRI{4,12},\rRI{13,19}}}
\end{tabproofwide}

\Needspace{14\baselineskip}
\FormulaThmDeltaK[Die Vereinigung bleibt eine Relation auf den gegebenen Mengen]{%
  \mathsf{WortStr}_A(W,e,s)\dsep x_0\in X\vdash G\subseteq W\times X
}{WordRecursionStageUnionSubset}{%
  \DeltaRow{Mengen}{A\dsep W\dsep e\dsep X\dsep x_0\dsep p\dsep n}\DeltaRow{Funktionen}{s\dsep r}}
\begin{tabproofwide}
  \proofstepwidestar[1]{\mathsf{WortStr}_A(W,e,s)}{\rA}
  \proofstepwidestar[2]{x_0\in X}{\rA}
  \proofstepwidestar[3]{p\in G}{\rA}
  \proofstepwidestar[1,2,3]{\exists n\in\mathbb N\;p\in H_n}{\rRE{3,\rLREa{\FormulaRefAuto{WordRecursionStageUnionMembership}[thm]{1,2}}}}
  \proofstepwidestar[5]{n\in\mathbb N\land p\in H_n}{\rA}
  \proofstepwidestar[1,2,5]{H_n\subseteq B}{\rAEa{\rRE{\rAEa{5},\rUE{\rAEb{\rAEb{\FormulaRefAuto{WordRecursionStagesEquations}[thm]{1,2}}}}}}}
  \proofstepwidestar[1,2,5]{p\in B}{\FormulaRefAuto{A\subseteq B\dsep x\in A\vdash x\in B}[thm]{6,\rAEb{5}}}
  \proofstepwidestar[1,2,3]{p\in B}{\rEE{4,5,7}}
  \proofstepwidestar[1,2]{G\subseteq B}{\FormulaRefAuto{SubsetIntroductionRule}[thm]{[3]\vdots8}}
  \proofstepwidestar[1,2]{G\subseteq W\times X}{\rIE{\FormulaRefAuto{WordRecursionAmbientDef}[def],9}}
\end{tabproofwide}

\Needspace{14\baselineskip}
\FormulaThmDeltaK[Die Vereinigung ist rechtseindeutig]{%
  \mathsf{WortStr}_A(W,e,s)\dsep x_0\in X\dsep r\colon X\times A\to X\vdash R(G)
}{WordRecursionStageUnionRightUnique}{%
  \DeltaRow{Mengen}{A\dsep W\dsep e\dsep X\dsep x_0\dsep t\dsep y\dsep z\dsep n\dsep m}\DeltaRow{Funktionen}{s\dsep r}}
\begin{tabproofwide}
  \proofstepwidestar[1]{\mathsf{WortStr}_A(W,e,s)}{\rA}
  \proofstepwidestar[2]{x_0\in X}{\rA}
  \proofstepwidestar[3]{r\colon X\times A\to X}{\rA}
  \proofstepwidestar[4]{(t,y)\in G\land(t,z)\in G}{\rA}
  \proofstepwidestar[1,2,4]{\exists n\in\mathbb N\;(t,y)\in H_n}{\rRE{\rAEa{4},\rLREa{\FormulaRefAuto{WordRecursionStageUnionMembership}[thm]{1,2}}}}
  \proofstepwidestar[1,2,4]{\exists m\in\mathbb N\;(t,z)\in H_m}{\rRE{\rAEb{4},\rLREa{\FormulaRefAuto{WordRecursionStageUnionMembership}[thm]{1,2}}}}
  \proofstepwidestar[7]{n\in\mathbb N\land(t,y)\in H_n}{\rA}
  \proofstepwidestar[8]{m\in\mathbb N\land(t,z)\in H_m}{\rA}
  \proofstepwidestar[7,8]{n\leq m\lor m\leq n}{\FormulaRefAuto{m\in\mathbb{N}\dsep n\in\mathbb{N}\vdash m\leq n\lor n\leq m}[thm]{\rAEa{7},\rAEa{8}}}
  \proofstepwidestar[10]{n\leq m}{\rA}
  \proofstepwidestar[1,2,3,7,8,10]{H_n\subseteq H_m}{\FormulaRefAuto{WordRecursionStagesChain}[thm]{1,2,3,\rAEa{7},\rAEa{8},10}}
  \proofstepwidestar[1,2,3,7,8,10]{(t,y)\in H_m}{\FormulaRefAuto{A\subseteq B\dsep x\in A\vdash x\in B}[thm]{11,\rAEb{7}}}
  \proofstepwidestar[1,2,3,8]{R(H_m)}{\rRE{\rAEa{8},\rUE{\FormulaRefAuto{WordRecursionStagesRightUnique}[thm]{1,2,3}}}}
  \proofstepwidestar[1,2,3,7,8,10]{y=z}{\rRE{\rAI{12,\rAEb{8}},\rUE{\rUE{\rUE{\FormulaRefAuto{WordRecursionRightUniqueDef}[def]{13}}}}}}
  \proofstepwidestar[15]{m\leq n}{\rA}
  \proofstepwidestar[1,2,3,7,8,15]{H_m\subseteq H_n}{\FormulaRefAuto{WordRecursionStagesChain}[thm]{1,2,3,\rAEa{8},\rAEa{7},15}}
  \proofstepwidestar[1,2,3,7,8,15]{(t,z)\in H_n}{\FormulaRefAuto{A\subseteq B\dsep x\in A\vdash x\in B}[thm]{16,\rAEb{8}}}
  \proofstepwidestar[1,2,3,7]{R(H_n)}{\rRE{\rAEa{7},\rUE{\FormulaRefAuto{WordRecursionStagesRightUnique}[thm]{1,2,3}}}}
  \proofstepwidestar[1,2,3,7,8,15]{y=z}{\rRE{\rAI{\rAEb{7},17},\rUE{\rUE{\rUE{\FormulaRefAuto{WordRecursionRightUniqueDef}[def]{18}}}}}}
  \proofstepwidestar[1,2,3,7,8]{y=z}{\rOE{9,10,14,15,19}}
  \proofstepwidestar[1,2,3,4,7]{y=z}{\rEE{6,8,20}}
  \proofstepwidestar[1,2,3,4]{y=z}{\rEE{5,7,21}}
  \proofstepwidestar[1,2,3]{R(G)}{\FormulaRefAuto{WordRecursionRightUniqueDef}[def]{\rUI{\rUI{\rUI{\rRI{4,22}}}}}}
\end{tabproofwide}

Für die Vollständigkeit wechseln wir nun von der Zahleninduktion zum
axiomatischen Wortinduktionsprinzip W3. Die Aussage lautet:
\emph{Dieses Wort wird in irgendeiner
Stufe erfasst.} Das Anfangswort wird in Stufe~\(0\) erfasst. Wird
\(u\) in Stufe~\(n\) erfasst, so wird jede Fortsetzung \(s(u,a)\)
in Stufe~\(\Succ(n)\) erfasst. Anders als bei einer einzelnen
Zahlenfolge kann es viele solche Fortsetzungen geben; der Schritt
verlangt die Aussage deshalb für jeden Buchstaben \(a\in A\).

\Needspace{14\baselineskip}
\FormulaThmDeltaK[Jedes Wort erhält in einer Stufe einen Wert]{%
  \begin{aligned}[t]
    &\mathsf{WortStr}_A(W,e,s)\dsep x_0\in X\dsep r\colon X\times A\to X\vdash{}\\[-2pt]
    &\forall u\in W\,\exists n\in\mathbb N\,\exists x\in X\;(u,x)\in H_n
  \end{aligned}
}{WordRecursionStagesTotal}{%
  \DeltaRow{Mengen}{A\dsep W\dsep e\dsep X\dsep x_0\dsep u\dsep v\dsep a\dsep n\dsep x}
  \DeltaRow{Funktionen}{s\dsep r}
  \DeltaRow{Abkürzung}{P(u)\coloneqq\exists n\in\mathbb N\,\exists x\in X\;(u,x)\in H_n}}
\begin{tabproofwide}
  \proofstepwidestar[1]{\mathsf{WortStr}_A(W,e,s)}{\rA}
  \proofstepwidestar[2]{x_0\in X}{\rA}
  \proofstepwidestar[3]{r\colon X\times A\to X}{\rA}
  \proofstepwidestar[1,2]{H_0=\{(e,x_0)\}\land\forall n\in\mathbb N\;(H_n\subseteq B\land H_{\Succ(n)}=T(H_n))}{\rAEb{\FormulaRefAuto{WordRecursionStagesEquations}[thm]{1,2}}}
  \proofstepwidestar[]{0\in\mathbb N}{\FormulaRefAuto{PeanoZero}[ax]}
  \proofstepwidestar[1,2]{(e,x_0)\in H_0}{\rIE{\FormulaRefAuto{a=b\vdash b=a}[thm]{\rAEa{4}},\FormulaRefAuto{a\in\{a\}}[thm]}}
  \proofstepwidestar[1,2]{P(e)}{\rEI{\rAI{5,\rEI{\rAI{2,6}}}}}
  \proofstepwidestar[8]{u\in W}{\rA}
  \proofstepwidestar[9]{a\in A}{\rA}
  \proofstepwidestar[10]{P(u)}{\rA}
  \proofstepwidestar[11]{n\in\mathbb N\land\exists x\in X\;(u,x)\in H_n}{\rA}
  \proofstepwidestar[12]{x\in X\land(u,x)\in H_n}{\rA}
  \proofstepwidestar[1,2,11]{H_n\subseteq B\land H_{\Succ(n)}=T(H_n)}{\rRE{\rAEa{11},\rUE{\rAEb{4}}}}
  \proofstepwidestar[1,2,3,11]{\forall v\in W\,\forall y\in X\,\forall b\in A\;((v,y)\in H_n\rightarrow(s(v,b),r(y,b))\in T(H_n))}{\rAEb{\FormulaRefAuto{WordRecursionStageBaseAndStep}[thm]{1,2,3,\rAEa{13}}}}
  \proofstepwidestar[1,2,3,8,9,11,12]{(s(u,a),r(x,a))\in T(H_n)}{\rRE{\rAEb{12},\rRE{9,\rUE{\rRE{\rAEa{12},\rUE{\rRE{8,\rUE{14}}}}}}}}
  \proofstepwidestar[1,2,3,8,9,11,12]{(s(u,a),r(x,a))\in H_{\Succ(n)}}{\rIE{\FormulaRefAuto{a=b\vdash b=a}[thm]{\rAEb{13}},15}}
  \proofstepwidestar[9,12]{(x,a)\in X\times A}{\FormulaRefAuto{(a,b)\in A\times B\eqvdash a\in A\land b\in B}[thm]{\rAI{\rAEa{12},9}}}
  \proofstepwidestar[3,9,12]{r(x,a)\in X}{\FormulaRefAuto{F\colon A\to B\dsep x\in A\vdash F(x)\in B}[thm]{3,17}}
  \proofstepwidestar[11]{\Succ(n)\in\mathbb N}{\FormulaRefAuto{PeanoSuccClosure}[ax]{\rAEa{11}}}
  \proofstepwidestar[1,2,3,8,9,11,12]{P(s(u,a))}{\rEI{\rAI{19,\rEI{\rAI{18,16}}}}}
  \proofstepwidestar[1,2,3,8,9,11]{P(s(u,a))}{\rEE{\rAEb{11},12,20}}
  \proofstepwidestar[1,2,3,8,9,10]{P(s(u,a))}{\rEE{10,11,21}}
  \proofstepwidestar[23]{v\in W}{\rA}
  \proofstepwidestar[1,2,3,23]{P(v)}{\FormulaRefAuto{WordStructureInductionRule}[ax]{1,7,[8,9,10]\vdots22,23}}
  \proofstepwidestar[1,2,3]{\forall u\in W\,\exists n\in\mathbb N\,\exists x\in X\;(u,x)\in H_n}{\rUI{\rRI{23,24}}}
\end{tabproofwide}

Nun sind beide benötigten Aussagen vorhanden: Jedes Wort besitzt einen
Wert, und zwei Werte desselben Worts sind gleich. Das Funktionskriterium
aus Band~5 macht die Vereinigung damit zu einer Funktion auf ganz \(W\).

\Needspace{14\baselineskip}
\FormulaThmDeltaK[Die Vereinigung der Stufen ist eine Funktion]{%
  \begin{aligned}[t]
    &\mathsf{WortStr}_A(W,e,s)\dsep x_0\in X\dsep r\colon X\times A\to X\vdash{}\\[-2pt]
    &G\colon W\to X
  \end{aligned}
}{WordStructureRecursionGraphFunction}{%
  \DeltaRow{Mengen}{A\dsep W\dsep e\dsep X\dsep x_0\dsep u\dsep n\dsep x\dsep y\dsep z}\DeltaRow{Funktionen}{s\dsep r}}
\begin{tabproofwide}
  \proofstepwidestar[1]{\mathsf{WortStr}_A(W,e,s)}{\rA}
  \proofstepwidestar[2]{x_0\in X}{\rA}
  \proofstepwidestar[3]{r\colon X\times A\to X}{\rA}
  \proofstepwidestar[1,2]{G\subseteq W\times X}{\FormulaRefAuto{WordRecursionStageUnionSubset}[thm]{1,2}}
  \proofstepwidestar[1,2,3]{R(G)}{\FormulaRefAuto{WordRecursionStageUnionRightUnique}[thm]{1,2,3}}
  \proofstepwidestar[6]{u\in W}{\rA}
  \proofstepwidestar[1,2,3,6]{\exists n\in\mathbb N\,\exists x\in X\;(u,x)\in H_n}{\rRE{6,\rUE{\FormulaRefAuto{WordRecursionStagesTotal}[thm]{1,2,3}}}}
  \proofstepwidestar[8]{n\in\mathbb N\land\exists x\in X\;(u,x)\in H_n}{\rA}
  \proofstepwidestar[9]{x\in X\land(u,x)\in H_n}{\rA}
  \proofstepwidestar[8,9]{\exists k\in\mathbb N\;(u,x)\in H_k}{\rEI{\rAI{\rAEa{8},\rAEb{9}}}}
  \proofstepwidestar[1,2,8,9]{(u,x)\in G}{\rRE{10,\rLREb{\FormulaRefAuto{WordRecursionStageUnionMembership}[thm]{1,2}}}}
  \proofstepwidestar[1,2,8,9]{\exists x\in X\;(u,x)\in G}{\rEI{\rAI{\rAEa{9},11}}}
  \proofstepwidestar[1,2,8]{\exists x\in X\;(u,x)\in G}{\rEE{\rAEb{8},9,12}}
  \proofstepwidestar[1,2,3,6]{\exists x\in X\;(u,x)\in G}{\rEE{7,8,13}}
  \proofstepwidestar[15]{y\in X\land(u,y)\in G}{\rA}
  \proofstepwidestar[16]{z\in X\land(u,z)\in G}{\rA}
  \proofstepwidestar[1,2,3,15,16]{y=z}{\rRE{\rAI{\rAEb{15},\rAEb{16}},\rUE{\rUE{\rUE{\FormulaRefAuto{WordRecursionRightUniqueDef}[def]{5}}}}}}
  \proofstepwidestar[1,2,3,6]{\exists!x\in X\;(u,x)\in G}{\UEI{14,15,16,17}}
  \proofstepwidestar[1,2,3]{\forall u\in W\,\exists!x\in X\;(u,x)\in G}{\rUI{\rRI{6,18}}}
  \proofstepwidestar[1,2,3]{G\colon W\to X}{\FormulaRefAuto{UniqueValuedGraphFunction}[thm]{4,19}}
\end{tabproofwide}

\subsubsection{Die Stufen liefern die Rekursionsgleichungen}

Zum Schluss lesen wir die Stufenkonstruktion als Funktionsgleichungen.
Das Anfangspaar liegt bereits in \(H_0\). Liegt \((u,x)\) in einer
Stufe, so liegt \((s(u,a),r(x,a))\) in der nächsten. Beide Paare
gehören deshalb zu \(G\). Dieser Übergang entspricht genau dem
Nachfolgerschritt der Zahlenrekursion; der zusätzliche Parameter \(a\)
hält fest, welche Fortsetzung des Worts ausgewertet wird.

\Needspace{14\baselineskip}
\FormulaThmDeltaK[Anfang und Fortsetzung liegen in der vereinigten Zuordnung]{%
  \begin{aligned}[t]
    &\mathsf{WortStr}_A(W,e,s)\dsep x_0\in X\dsep r\colon X\times A\to X\vdash{}\\[-2pt]
    &(e,x_0)\in G\land\forall u\in W\,\forall x\in X\,\forall a\in A\;\\[-2pt]
    &\qquad((u,x)\in G\rightarrow(s(u,a),r(x,a))\in G)
  \end{aligned}
}{WordRecursionStageUnionBaseAndStep}{%
  \DeltaRow{Mengen}{A\dsep W\dsep e\dsep X\dsep x_0\dsep u\dsep x\dsep a\dsep n}\DeltaRow{Funktionen}{s\dsep r}}
\begin{tabproofwide}
  \proofstepwidestar[1]{\mathsf{WortStr}_A(W,e,s)}{\rA}
  \proofstepwidestar[2]{x_0\in X}{\rA}
  \proofstepwidestar[3]{r\colon X\times A\to X}{\rA}
  \proofstepwidestar[1,2]{H_0=\{(e,x_0)\}\land\forall n\in\mathbb N\;(H_n\subseteq B\land H_{\Succ(n)}=T(H_n))}{\rAEb{\FormulaRefAuto{WordRecursionStagesEquations}[thm]{1,2}}}
  \proofstepwidestar[]{0\in\mathbb N}{\FormulaRefAuto{PeanoZero}[ax]}
  \proofstepwidestar[1,2]{(e,x_0)\in H_0}{\rIE{\FormulaRefAuto{a=b\vdash b=a}[thm]{\rAEa{4}},\FormulaRefAuto{a\in\{a\}}[thm]}}
  \proofstepwidestar[1,2]{(e,x_0)\in G}{\rRE{\rLREb{\FormulaRefAuto{WordRecursionStageUnionMembership}[thm]{1,2}},\rEI{\rAI{5,6}}}}
  \proofstepwidestar[8]{u\in W}{\rA}
  \proofstepwidestar[9]{x\in X}{\rA}
  \proofstepwidestar[10]{a\in A}{\rA}
  \proofstepwidestar[11]{(u,x)\in G}{\rA}
  \proofstepwidestar[1,2,11]{\exists n\in\mathbb N\;(u,x)\in H_n}{\rRE{11,\rLREa{\FormulaRefAuto{WordRecursionStageUnionMembership}[thm]{1,2}}}}
  \proofstepwidestar[13]{n\in\mathbb N\land(u,x)\in H_n}{\rA}
  \proofstepwidestar[1,2,13]{H_n\subseteq B\land H_{\Succ(n)}=T(H_n)}{\rRE{\rAEa{13},\rUE{\rAEb{4}}}}
  \proofstepwidestar[1,2,3,13]{\forall v\in W\,\forall y\in X\,\forall b\in A\;((v,y)\in H_n\rightarrow(s(v,b),r(y,b))\in T(H_n))}{\rAEb{\FormulaRefAuto{WordRecursionStageBaseAndStep}[thm]{1,2,3,\rAEa{14}}}}
  \proofstepwidestar[1,2,3,8,9,10,13]{(s(u,a),r(x,a))\in T(H_n)}{\rRE{\rAEb{13},\rRE{10,\rUE{\rRE{9,\rUE{\rRE{8,\rUE{15}}}}}}}}
  \proofstepwidestar[1,2,3,8,9,10,13]{(s(u,a),r(x,a))\in H_{\Succ(n)}}{\rIE{\FormulaRefAuto{a=b\vdash b=a}[thm]{\rAEb{14}},16}}
  \proofstepwidestar[13]{\Succ(n)\in\mathbb N}{\FormulaRefAuto{PeanoSuccClosure}[ax]{\rAEa{13}}}
  \proofstepwidestar[1,2,3,8,9,10,13]{\exists k\in\mathbb N\;(s(u,a),r(x,a))\in H_k}{\rEI{\rAI{18,17}}}
  \proofstepwidestar[1,2,3,8,9,10,13]{(s(u,a),r(x,a))\in G}{\rRE{19,\rLREb{\FormulaRefAuto{WordRecursionStageUnionMembership}[thm]{1,2}}}}
  \proofstepwidestar[1,2,3,8,9,10,11]{(s(u,a),r(x,a))\in G}{\rEE{12,13,20}}
  \proofstepwidestar[1,2,3]{\forall u\in W\,\forall x\in X\,\forall a\in A\;((u,x)\in G\rightarrow(s(u,a),r(x,a))\in G)}{\rUI{\rRI{8,\rUI{\rRI{9,\rUI{\rRI{10,\rRI{11,21}}}}}}}}
  \proofstepwidestar[1,2,3]{\begin{aligned}[t]&(e,x_0)\in G\land{}\\[-2pt]&\forall u\in W\,\forall x\in X\,\forall a\in A\;\\[-2pt]&\quad((u,x)\in G\rightarrow(s(u,a),r(x,a))\in G)\end{aligned}}{\rAI{7,22}}
\end{tabproofwide}

\Needspace{14\baselineskip}
\FormulaThmDeltaK[Die vereinigte Funktion erfüllt die Wortrekursion]{%
  \begin{aligned}[t]
    &\mathsf{WortStr}_A(W,e,s)\dsep x_0\in X\dsep r\colon X\times A\to X\vdash{}\\[-2pt]
    &\mathsf{WRec}_A(G;W,e,s;X,x_0,r)
  \end{aligned}
}{WordStructureRecursionGraphEquations}{%
  \DeltaRow{Mengen}{A\dsep W\dsep e\dsep X\dsep x_0\dsep u\dsep a}\DeltaRow{Funktionen}{s\dsep r}}
\begin{tabproofwide}
  \proofstepwidestar[1]{\mathsf{WortStr}_A(W,e,s)}{\rA}
  \proofstepwidestar[2]{x_0\in X}{\rA}
  \proofstepwidestar[3]{r\colon X\times A\to X}{\rA}
  \proofstepwidestar[1,2,3]{G\colon W\to X}{\FormulaRefAuto{WordStructureRecursionGraphFunction}[thm]{1,2,3}}
  \proofstepwidestar[1,2,3]{\begin{aligned}[t]&(e,x_0)\in G\land{}\\[-2pt]&\forall u\in W\,\forall x\in X\,\forall a\in A\;\\[-2pt]&\quad((u,x)\in G\rightarrow(s(u,a),r(x,a))\in G)\end{aligned}}{\FormulaRefAuto{WordRecursionStageUnionBaseAndStep}[thm]{1,2,3}}
  \proofstepwidestar[1,2,3]{G(e)=x_0}{\FormulaRefAuto{F\colon A\to B\dsep(x,y)\in F\vdash F(x)=y}[thm]{4,\rAEa{5}}}
  \proofstepwidestar[7]{u\in W}{\rA}
  \proofstepwidestar[8]{a\in A}{\rA}
  \proofstepwidestar[1,2,3,7]{G(u)\in X}{\FormulaRefAuto{F\colon A\to B\dsep x\in A\vdash F(x)\in B}[thm]{4,7}}
  \proofstepwidestar[1,2,3,7]{(u,G(u))\in G}{\FormulaRefAuto{F\colon A\to B\dsep x\in A\vdash(x,F(x))\in F}[thm]{4,7}}
  \proofstepwidestar[1,2,3,7,8]{(s(u,a),r(G(u),a))\in G}{\rRE{10,\rRE{8,\rUE{\rRE{9,\rUE{\rRE{7,\rUE{\rAEb{5}}}}}}}}}
  \proofstepwidestar[1,2,3,7,8]{G(s(u,a))=r(G(u),a)}{\FormulaRefAuto{F\colon A\to B\dsep(x,y)\in F\vdash F(x)=y}[thm]{4,11}}
  \proofstepwidestar[1,2,3]{\forall u\in W\,\forall a\in A\;G(s(u,a))=r(G(u),a)}{\rUI{\rRI{7,\rUI{\rRI{8,12}}}}}
  \proofstepwidestar[1,2,3]{\mathsf{WRec}_A(G;W,e,s;X,x_0,r)}{\FormulaRefAuto{WordRecursionSpecificationDef}[def]{\rAI{4,\rAI{6,13}}}}
\end{tabproofwide}

Damit ist die Existenz bewiesen: Die Rekursion auf \(\mathbb N\)
konstruiert die Stufen, die Wortaxiome sichern ihre widerspruchsfreie
Fortsetzung, und die Wortinduktion zeigt, dass kein Wort ausgelassen
wird. Zusammen mit der vorher bewiesenen Eindeutigkeit ergibt dies
den angekündigten Wortrekursionssatz.
